\section{Apply one complete record}
The XOR kernel has no knowledge of a file format. Give the next routine,
\texttt{video\_apply\_record}, responsibility for one encoded frame record. Its input
starts at the record's kind byte, includes the sixteen-byte header, and ends at the
last payload byte. The global file header is outside this window. A later file
reader will supply that exact window after locating the record.

Pass the record address in \texttt{\$a0}, its byte count in \texttt{\$a1}, the working
image address in \texttt{\$a2}, and a context address in \texttt{\$a3}. The record
may start at an unaligned byte address; the image must be aligned to four bytes and
the context to eight. The caller supplies readable record bytes, writable image
storage of at least 16\,384 bytes, a writable 48-byte context, and a valid private
stack. These windows are separate and do not overlap. Capacity is a caller's
promise about storage; the routine cannot discover whether an arbitrary numeric
address names a real allocation.

\begin{table}[htbp]
\centering
\caption{Native context for a single record call. Each field is eight bytes.}
\label{tab:video-record-context}
\begin{tabular}{r p{1.4in} p{2.9in}}
\toprule
Offset & Field & Meaning \\
\midrule
0 & Frame capacity & Input, at least 16\,384 bytes. \\
8 & Expected sequence & Input from 0 through 255; incremented on success. \\
16 & Full-frame period & Input from 1 through 64. \\
24 & Last kind & Output, zero for full or one for delta. \\
32 & Active frame bytes & Output, 16\,384 on success. \\
40 & Retained suffix & Unused by this routine; leave its bytes unchanged. \\
\bottomrule
\end{tabular}
\end{table}

The context uses the simulator's selected native byte order. Record integers keep
their specified big-endian order in either simulator setting. For a delta, the
working image contains the previously accepted frame as 4096 numeric
\texttt{0x00RRGGBB} words. A full frame replaces that image and does not depend on
its earlier contents.

Return zero in \texttt{\$v0} and 16\,384 in \texttt{\$v1} on success. Only then
increment the expected sequence and publish the kind and active byte count. Return
one and zero on failure, retaining every image and context byte. The source remains
immutable in both cases. Preserve saved registers, the frame pointer and the
caller's stack pointer. The final accepted sequence may be 255: publication advances
the context to 256, and a subsequent call is rejected. File-level frame-count and
end-of-file checks will belong to the reader, which knows the complete file.

The supplied caller is in \texttt{packages/034-lossless-video/construction/video-record-construction/main.asm}. Save
a working copy and fill only the marked regions. The small
\texttt{record.asm} file contains the same unfinished routine for inspection;
\texttt{main.asm} also supplies the data and executable caller. Do not concatenate
the two routines. Prepare an observation with, for example,
\texttt{python3 packages/034-lossless-video/construction/video-record-construction/prepare.py s07-1 --source work.asm
--output observe.asm}. The output must be a new file. Preparation replaces fixture
data while retaining your instruction text.

An unfinished region returns development status four and count zero. This is a
construction aid. After either application loop has been written, it can change the
image before the unfinished publication region returns four. A completed routine
has only success zero or failure one; failure one never permits partial application.
Stop at the named observation before continuing. Reaching a checkpoint shows what
has been established at that boundary; it is not a substitute for the complete
success and failure tests.
